% ============================================================================
% figures.tex
% ----------------------------------------------------------------------------
% Required TikZ libraries (in main preamble):
% positioning, arrows.meta, calc, shapes.geometric, fit, backgrounds, matrix
%
% This file defines reusable figure macros for insertion at section-level
% locations in the paper. Each macro expands to a full figure environment.
% ============================================================================

% Unified figure colors (formal, muted)
\definecolor{shNeutral}{RGB}{45,45,45}
\definecolor{shBlue}{RGB}{44,90,140}
\definecolor{shRed}{RGB}{130,48,48}
\definecolor{shGreen}{RGB}{72,120,78}

% Unified drawing styles
\tikzset{
  sh/edge/.style={draw=shNeutral, line width=0.95pt},
  sh/edgeBlue/.style={draw=shBlue, line width=1.0pt},
  sh/edgeRed/.style={draw=shRed, line width=1.0pt},
  sh/couple/.style={draw=shRed, line width=0.95pt, dashed},
  sh/infer/.style={draw=shNeutral, line width=0.9pt, dotted},
  sh/node/.style={circle, draw=shNeutral, line width=0.9pt, fill=white, inner sep=1.45pt},
  sh/nodeBlue/.style={circle, draw=shBlue, line width=0.95pt, fill=white, inner sep=1.45pt},
  sh/nodeGreen/.style={circle, draw=shGreen, line width=0.95pt, fill=white, inner sep=1.45pt},
  sh/arrow/.style={-{Stealth[length=2.2mm,width=1.6mm]}, line width=0.95pt, draw=shNeutral},
  sh/arrowBlue/.style={-{Stealth[length=2.2mm,width=1.6mm]}, line width=0.95pt, draw=shBlue},
  sh/arrowRed/.style={-{Stealth[length=2.2mm,width=1.6mm]}, line width=0.95pt, draw=shRed},
  sh/box/.style={draw=shNeutral, rounded corners=1.2pt, line width=0.9pt, fill=white},
  sh/boxBlue/.style={draw=shBlue, rounded corners=1.2pt, line width=0.9pt, fill=shBlue!7},
  sh/boxRed/.style={draw=shRed, rounded corners=1.2pt, line width=0.9pt, fill=shRed!8},
  sh/boxGreen/.style={draw=shGreen, rounded corners=1.2pt, line width=0.9pt, fill=shGreen!8}
}

% ============================================================================
% Figure A1
% Name: Progressive construction to Supra-Hodge
% Purpose: Show left-to-right construction pipeline from graph to coupled
%          higher-order multi-layer system.
% Insert near: Introduction (replaces old conceptual hierarchy figure)
% ============================================================================
\newcommand{\FigureAOne}{
\begin{figure*}[t]
\centering
\begin{tikzpicture}[x=1cm,y=1cm,>=Latex,scale=1.04,every node/.style={transform shape}]

% ----------------------------
% Color palette
% ----------------------------
\definecolor{topoblue}{RGB}{70,120,200}
\definecolor{couplered}{RGB}{190,95,95}
\definecolor{harmgreen}{RGB}{80,150,110}
\definecolor{ghostgray}{RGB}{180,180,180}
\definecolor{framegray}{RGB}{80,80,80}
\definecolor{lightfill}{RGB}{245,247,250}

% ----------------------------
% Shared styles
% ----------------------------
\tikzset{
    baseNode/.style={
        circle, draw=framegray, fill=white,
        inner sep=0pt, minimum size=4.6pt, line width=0.8pt
    },
    phasePanel/.style={rounded corners=4pt, draw=ghostgray!85, fill=lightfill, line width=0.7pt},
    neutralEdge/.style={draw=framegray, line width=0.9pt},
    topoEdge/.style={draw=topoblue, line width=1.0pt},
    couplingEdge/.style={draw=couplered!75, dashed, line width=0.8pt},
    ghostEdge/.style={draw=ghostgray, dashed, line width=0.7pt},
    panelTitle/.style={font=\small},
    miniLabel/.style={font=\scriptsize},
    numLabel/.style={font=\scriptsize, text=framegray},
    morphArrow/.style={->, draw=framegray, line width=0.9pt},
    lpbox/.style={rounded corners=4pt, draw=framegray, fill=lightfill, line width=0.8pt}
}

% ----------------------------
% Panel centers
% ----------------------------
\coordinate (P1) at (1.8,2.0);
\coordinate (P2) at (5.8,2.0);
\coordinate (P3) at (10.3,2.0);
\coordinate (P4) at (15.2,2.0);

% ----------------------------
% Panel 1: Graph
% ----------------------------
\draw[phasePanel] ($(P1)+(-1.35,-1.20)$) rectangle ($(P1)+(1.35,1.25)$);
\coordinate (p1a) at ($(P1)+(-0.85,-0.85)$);
\coordinate (p1b) at ($(P1)+( 0.85,-0.85)$);
\coordinate (p1c) at ($(P1)+(-0.85, 0.85)$);
\coordinate (p1d) at ($(P1)+( 0.85, 0.85)$);

\draw[neutralEdge] (p1a)--(p1b)--(p1d)--(p1c)--cycle;
\draw[neutralEdge] (p1a)--(p1d);

\node[baseNode] at (p1a) {};
\node[baseNode] at (p1b) {};
\node[baseNode] at (p1c) {};
\node[baseNode] at (p1d) {};

\node[numLabel, below=2pt of p1a] {1};
\node[numLabel, below=2pt of p1b] {2};
\node[numLabel, left=2pt of p1c] {3};
\node[numLabel, right=2pt of p1d] {4};

\node[panelTitle] at ($(P1)+(0,1.70)$) {(a) Graph};
\node[miniLabel]  at ($(P1)+(0,-1.72)$) {pairwise skeleton};

% ----------------------------
% Panel 2: Simplicial complex
% ----------------------------
\draw[phasePanel] ($(P2)+(-1.35,-1.20)$) rectangle ($(P2)+(1.35,1.25)$);
\coordinate (p2a) at ($(P2)+(-0.85,-0.85)$);
\coordinate (p2b) at ($(P2)+( 0.85,-0.85)$);
\coordinate (p2c) at ($(P2)+(-0.85, 0.85)$);
\coordinate (p2d) at ($(P2)+( 0.85, 0.85)$);

% filled triangle (1,2,4)
\fill[topoblue!20] (p2a)--(p2b)--(p2d)--cycle;
% ghost triangle (1,3,4)
\fill[topoblue!8] (p2a)--(p2c)--(p2d)--cycle;

\draw[neutralEdge] (p2a)--(p2b)--(p2d)--(p2c)--cycle;
\draw[neutralEdge] (p2a)--(p2d);

\node[baseNode] at (p2a) {};
\node[baseNode] at (p2b) {};
\node[baseNode] at (p2c) {};
\node[baseNode] at (p2d) {};

\node[numLabel, below=2pt of p2a] {1};
\node[numLabel, below=2pt of p2b] {2};
\node[numLabel, left=2pt of p2c] {3};
\node[numLabel, right=2pt of p2d] {4};

\node[panelTitle] at ($(P2)+(0,1.70)$) {(b) Simplicial complex};
\node[miniLabel]  at ($(P2)+(0,-1.72)$) {higher-order closure};

% ----------------------------
% Panel 3: Multilayer graph
% ----------------------------
\draw[phasePanel] ($(P3)+(-1.55,-1.95)$) rectangle ($(P3)+(1.35,1.45)$);
% top layer
\coordinate (p3ta) at ($(P3)+(-0.70,-0.30)$);
\coordinate (p3tb) at ($(P3)+( 0.70,-0.30)$);
\coordinate (p3tc) at ($(P3)+(-0.70, 1.10)$);
\coordinate (p3td) at ($(P3)+( 0.70, 1.10)$);

% bottom layer
\coordinate (p3ba) at ($(P3)+(-0.70,-1.70)$);
\coordinate (p3bb) at ($(P3)+( 0.70,-1.70)$);
\coordinate (p3bc) at ($(P3)+(-0.70,-0.30)$);
\coordinate (p3bd) at ($(P3)+( 0.70,-0.30)$);

\draw[neutralEdge] (p3ta)--(p3tb)--(p3td)--(p3tc)--cycle;
\draw[neutralEdge] (p3ta)--(p3td);

\draw[neutralEdge] (p3ba)--(p3bb)--(p3bd)--(p3bc)--cycle;
\draw[neutralEdge] (p3ba)--(p3bd);

\draw[ghostEdge] (p3ta)--(p3ba);
\draw[ghostEdge] (p3tb)--(p3bb);
\draw[ghostEdge] (p3tc)--(p3bc);
\draw[ghostEdge] (p3td)--(p3bd);

\node[baseNode] at (p3ta) {};
\node[baseNode] at (p3tb) {};
\node[baseNode] at (p3tc) {};
\node[baseNode] at (p3td) {};
\node[baseNode] at (p3ba) {};
\node[baseNode] at (p3bb) {};
\node[baseNode] at (p3bc) {};
\node[baseNode] at (p3bd) {};

\node[miniLabel, anchor=west] at ($(P3)+(-1.42,0.48)$) {Layer A};
\node[miniLabel, anchor=west] at ($(P3)+(-1.42,-1.05)$) {Layer B};

\node[panelTitle] at ($(P3)+(0,1.95)$) {(c) Multilayer graph};
\node[miniLabel]  at ($(P3)+(0,-2.20)$) {same entities, multiple views};

% ----------------------------
% Panel 4: Supra-Hodge setting
% ----------------------------
% outer operator box
\draw[lpbox] ($(P4)+(-1.70,-2.00)$) rectangle ($(P4)+(1.70,1.70)$);
\node[font=\small] at ($(P4)+(0,1.78)$) {$L_p$};

% top layer
\coordinate (p4ta) at ($(P4)+(-0.95,-0.25)$);
\coordinate (p4tb) at ($(P4)+( 0.20,-0.25)$);
\coordinate (p4tc) at ($(P4)+(-0.95, 0.95)$);
\coordinate (p4td) at ($(P4)+( 0.20, 0.95)$);

% bottom layer
\coordinate (p4ba) at ($(P4)+(-0.20,-1.35)$);
\coordinate (p4bb) at ($(P4)+( 0.95,-1.35)$);
\coordinate (p4bc) at ($(P4)+(-0.20,-0.15)$);
\coordinate (p4bd) at ($(P4)+( 0.95,-0.15)$);

% higher-order faces
\fill[topoblue!20] (p4ta)--(p4tb)--(p4td)--cycle;
\fill[topoblue!20] (p4ba)--(p4bc)--(p4bd)--cycle;

% intra-layer edges
\draw[topoEdge] (p4ta)--(p4tb)--(p4td)--(p4tc)--cycle;
\draw[topoEdge] (p4ta)--(p4td);

\draw[topoEdge] (p4ba)--(p4bb)--(p4bd)--(p4bc)--cycle;
\draw[topoEdge] (p4ba)--(p4bd);

% inter-layer couplings
\draw[couplingEdge] (p4ta)--(p4ba);
\draw[couplingEdge] (p4tb)--(p4bb);
\draw[couplingEdge] (p4tc)--(p4bc);
\draw[couplingEdge] (p4td)--(p4bd);

\node[baseNode] at (p4ta) {};
\node[baseNode] at (p4tb) {};
\node[baseNode] at (p4tc) {};
\node[baseNode] at (p4td) {};
\node[baseNode] at (p4ba) {};
\node[baseNode] at (p4bb) {};
\node[baseNode] at (p4bc) {};
\node[baseNode] at (p4bd) {};

\node[panelTitle] at ($(P4)+(0,2.28)$) {(d) Supra-Hodge setting};

% ----------------------------
% Morph arrows between panels
% ----------------------------
\draw[morphArrow] ($(P1)+(1.35,0)$) to[out=0,in=180] ($(P2)+(-1.35,0)$);
\draw[morphArrow] ($(P2)+(1.35,0)$) to[out=0,in=180] ($(P3)+(-1.35,0)$);
\draw[morphArrow] ($(P3)+(1.45,0)$) to[out=0,in=180] ($(P4)+(-2.10,0)$);

\end{tikzpicture}
\caption{
Progressive emergence of the Supra-Hodge setting: a pairwise graph becomes a simplicial complex by adding higher-order faces, then a multilayer graph by duplicating the entity set across views, and finally a coupled higher-order system in which intra-layer topology and inter-layer coupling coexist inside a single operator model.
}
\label{fig:construction}
\end{figure*}
}

% ============================================================================
% Figure A2
% Name: Chain spaces and signals
% Purpose: Visualize C_0, C_1, C_2 signal types on a common simplicial base.
% Insert near: Simplicial/chain-space motivation text
% ============================================================================
\newcommand{\FigureATwo}{
\begin{figure*}[t]
\centering
\begin{tikzpicture}[x=1cm,y=1cm,>=Latex,scale=1.10,every node/.style={transform shape}]

% ----------------------------
% Color palette
% ----------------------------
\definecolor{topoblue}{RGB}{70,120,200}
\definecolor{couplered}{RGB}{190,95,95}
\definecolor{harmgreen}{RGB}{80,150,110}
\definecolor{softgold}{RGB}{210,170,80}
\definecolor{ghostgray}{RGB}{185,185,185}
\definecolor{framegray}{RGB}{80,80,80}
\definecolor{lightfill}{RGB}{246,247,249}

% ----------------------------
% Styles
% ----------------------------
\tikzset{
    baseNode/.style={
        circle, draw=framegray, fill=white,
        inner sep=0pt, minimum size=5.2pt, line width=0.8pt
    },
    paleNode/.style={
        circle, draw=framegray!80, fill=white!95!ghostgray,
        inner sep=0pt, minimum size=5.0pt, line width=0.7pt
    },
    framePanel/.style={
        rounded corners=3pt, draw=framegray, line width=0.75pt
    },
    faintEdge/.style={draw=ghostgray, line width=0.7pt},
    mainEdge/.style={draw=framegray, line width=0.9pt},
    flowArrow/.style={->, draw=topoblue, line width=1.0pt},
    legendArrow/.style={->, draw=topoblue, line width=0.9pt},
    panelTitle/.style={font=\small},
    miniLabel/.style={font=\scriptsize},
    noteLabel/.style={font=\scriptsize, text=framegray},
    valueLabel/.style={font=\scriptsize},
    legendDot/.style={circle, draw=framegray, minimum size=5pt, inner sep=0pt}
}

% ----------------------------
% Panel frames
% ----------------------------
\draw[framePanel] (0.4,0.6) rectangle (4.2,4.2);
\draw[framePanel] (4.9,0.6) rectangle (8.7,4.2);
\draw[framePanel] (9.4,0.6) rectangle (13.2,4.2);

% ----------------------------
% Shared complex coordinates
% ----------------------------
% Left panel offset = (0.4,0.6)
\coordinate (LA) at (1.2,1.6);
\coordinate (LB) at (2.4,1.2);
\coordinate (LC) at (3.6,1.8);
\coordinate (LD) at (2.9,3.0);
\coordinate (LE) at (1.6,3.3);

% Mid panel offset = +4.5 in x
\coordinate (MA) at ($(LA)+(4.5,0)$);
\coordinate (MB) at ($(LB)+(4.5,0)$);
\coordinate (MC) at ($(LC)+(4.5,0)$);
\coordinate (MD) at ($(LD)+(4.5,0)$);
\coordinate (ME) at ($(LE)+(4.5,0)$);

% Right panel offset = +9.0 in x
\coordinate (RA) at ($(LA)+(9.0,0)$);
\coordinate (RB) at ($(LB)+(9.0,0)$);
\coordinate (RC) at ($(LC)+(9.0,0)$);
\coordinate (RD) at ($(LD)+(9.0,0)$);
\coordinate (RE) at ($(LE)+(9.0,0)$);

% ----------------------------
% Panel 1: Node signals C_0
% ----------------------------
\draw[faintEdge] (LA)--(LB)--(LC);
\draw[faintEdge] (LA)--(LE)--(LD)--(LC);
\draw[faintEdge] (LB)--(LD);
\draw[faintEdge] (LA)--(LD);

% Node fills by value
\fill[topoblue!35]  (LA) circle (0.13);
\fill[harmgreen!28] (LB) circle (0.13);
\fill[softgold!45]  (LC) circle (0.13);
\fill[ghostgray!70] (LD) circle (0.13);
\fill[couplered!35] (LE) circle (0.13);

\node[baseNode] at (LA) {};
\node[baseNode] at (LB) {};
\node[baseNode] at (LC) {};
\node[baseNode] at (LD) {};
\node[baseNode] at (LE) {};

% value labels
\node[valueLabel, text=topoblue!90!black]   at ($(LA)+(-0.18,0.28)$) {$0.9$};
\node[valueLabel, text=harmgreen!70!black]  at ($(LB)+(-0.20,-0.28)$) {$-0.4$};
\node[valueLabel, text=softgold!60!black]   at ($(LC)+(0.24,0.12)$) {$0.6$};
\node[valueLabel, text=framegray]           at ($(LD)+(0.18,0.24)$) {$0.2$};

\node[panelTitle] at (2.3,4.45) {Node signals $C_0$};

% ----------------------------
% Panel 2: Edge signals C_1
% ----------------------------
\draw[faintEdge] (MA)--(MB)--(MC);
\draw[faintEdge] (MA)--(ME)--(MD)--(MC);
\draw[faintEdge] (MB)--(MD);
\draw[faintEdge] (MA)--(MD);

\node[paleNode] at (MA) {};
\node[paleNode] at (MB) {};
\node[paleNode] at (MC) {};
\node[paleNode] at (MD) {};
\node[paleNode] at (ME) {};

% oriented flows
\draw[flowArrow, shorten <=2pt, shorten >=2pt] (MA)--(MB);
\draw[flowArrow, shorten <=2pt, shorten >=2pt, line width=1.15pt] (MA)--(MD);
\draw[flowArrow, shorten <=2pt, shorten >=2pt, line width=0.95pt] (MD)--(MB);
\draw[flowArrow, shorten <=2pt, shorten >=2pt] (MC)--(MD);
\draw[flowArrow, shorten <=2pt, shorten >=2pt, line width=0.8pt] (ME)--(MA);

\node[valueLabel, anchor=west] at ($(MB)+(0.28,-0.36)$) {$f \in C_1$};

\node[panelTitle] at (6.8,4.55) {Edge signals $C_1$};

% ----------------------------
% Panel 3: Triangle signals C_2
% ----------------------------
% faces
\fill[topoblue!22] (RA)--(RB)--(RD)--cycle;
\fill[topoblue!12] (RA)--(RD)--(RE)--cycle;

\draw[faintEdge] (RA)--(RB)--(RC);
\draw[faintEdge] (RA)--(RE)--(RD)--(RC);
\draw[faintEdge] (RB)--(RD);
\draw[faintEdge] (RA)--(RD);

\node[paleNode] at (RA) {};
\node[paleNode] at (RB) {};
\node[paleNode] at (RC) {};
\node[paleNode] at (RD) {};
\node[paleNode] at (RE) {};

\node[valueLabel, anchor=west, fill=white, inner sep=1pt] at ($(RB)+(0.28,-0.50)$) {$g \in C_2$};

% small orientation marker
\draw[topoblue, line width=0.8pt, ->] ($(RB)+(0.04,-0.32)$) arc[start angle=45,end angle=330,radius=0.23];

\node[panelTitle] at (11.3,4.55) {Triangle signals $C_2$};

% ----------------------------
% Bottom subtitles
% ----------------------------
\node[miniLabel] at (2.3,0.18) {Pointwise state};
\node[miniLabel] at (6.8,0.18) {Relational flow};
\node[miniLabel] at (11.3,0.18) {Joint interaction};

% ----------------------------
% Legend strip
% ----------------------------
\draw[rounded corners=8pt, draw=ghostgray!70, fill=lightfill] (0.85,-0.82) rectangle (12.85,-0.02);

% scalar amplitude
\fill[topoblue!25] (1.6,-0.35) circle (0.12);
\node[baseNode] at (1.6,-0.35) {};
\node[noteLabel, anchor=west] at (1.90,-0.42) {scalar amplitude};

% orientation
\draw[legendArrow] (4.45,-0.42) -- (5.45,-0.42);
\node[noteLabel, anchor=west] at (5.70,-0.42) {oriented edge signal};

% simplex intensity
\fill[topoblue!20] (9.00,-0.54) -- (9.34,-0.42) -- (9.00,-0.30) -- (8.66,-0.42) -- cycle;
\node[noteLabel, anchor=west] at (9.58,-0.42) {simplex intensity / circulation};

\end{tikzpicture}
\caption{
Signals on simplices of different order over the same underlying simplicial complex: node signals in $C_0$, oriented edge signals in $C_1$, and triangle-level signals in $C_2$. The display treats the three signal types as parallel views of the same base higher-order structure.
}
\label{fig:chains}
\end{figure*}
}

% ============================================================================
% Figure A3
% Name: Alignment into ambient chain space
% Purpose: Formal embedding diagram for layerwise chain spaces into common Ctil_p.
% Insert near: Coupling/alignment section after embedding definition
% ============================================================================
\newcommand{\FigureAThree}{
\begin{figure*}[t]
\centering
\begin{tikzpicture}[x=1cm,y=1cm,>=Latex,scale=1.12,every node/.style={transform shape}]

% ----------------------------
% Color palette
% ----------------------------
\definecolor{topoblue}{RGB}{70,120,200}
\definecolor{couplered}{RGB}{190,95,95}
\definecolor{harmgreen}{RGB}{80,150,110}
\definecolor{ghostgray}{RGB}{180,180,180}
\definecolor{framegray}{RGB}{80,80,80}
\definecolor{lightfill}{RGB}{245,247,250}

% ----------------------------
% Styles
% ----------------------------
\tikzset{
    sheet/.style={
        rounded corners=3pt,
        draw=framegray,
        fill=lightfill,
        line width=0.8pt
    },
    ambient/.style={
        rounded corners=3pt,
        draw=framegray,
        fill=white,
        line width=0.85pt
    },
    gate/.style={
        rounded corners=2pt,
        draw=framegray,
        fill=lightfill,
        line width=0.8pt
    },
    simplexNode/.style={
        circle, draw=framegray, fill=white,
        inner sep=0pt, minimum size=4.4pt, line width=0.7pt
    },
    simplexEdge/.style={draw=topoblue, line width=0.9pt},
    faintEdge/.style={draw=ghostgray, line width=0.7pt},
    mapArrow/.style={->, draw=framegray, line width=0.9pt},
    smallLabel/.style={font=\scriptsize},
    mainLabel/.style={font=\small},
    coordLabel/.style={font=\scriptsize, text=framegray},
    zeroLabel/.style={font=\scriptsize, text=ghostgray},
    solidSlot/.style={draw=topoblue!70, line width=0.7pt},
    zeroSlot/.style={draw=ghostgray!75, dashed, line width=0.6pt}
}

% ----------------------------
% Left source sheets
% ----------------------------

% Sheet 1 center
\coordinate (S1) at (2.3,3.45);
% Sheet 2 center
\coordinate (S2) at (2.3,1.45);

% Sheet 1 rotated box
\begin{scope}[shift={(S1)}, rotate=0]
    \draw[sheet] (-1.5,-0.65) rectangle (1.5,0.65);
    \node[mainLabel] at (0,0.95) {$C_p\!\left(K^{(1)}\right)$};

    % miniature simplex icons / coordinates
    \coordinate (s1a) at (-0.95,-0.05);
    \coordinate (s1b) at (-0.45,0.28);
    \coordinate (s1c) at (0.10,-0.02);
    \coordinate (s1d) at (0.65,0.24);

    \draw[simplexEdge] (s1a)--(s1b)--(s1c);
    \draw[simplexEdge] (s1a)--(s1c);
    \fill[topoblue!18] (s1a)--(s1b)--(s1c)--cycle;
    \draw[simplexEdge] (s1c)--(s1d);

    \node[simplexNode] at (s1a) {};
    \node[simplexNode] at (s1b) {};
    \node[simplexNode] at (s1c) {};
    \node[simplexNode] at (s1d) {};

    \node[coordLabel] at (-0.95,-0.42) {$e_1$};
    \node[coordLabel] at (-0.12,-0.42) {$e_2$};
    \node[coordLabel] at (0.72,-0.42) {$e_4$};
\end{scope}

% Sheet 2 rotated box
\begin{scope}[shift={(S2)}, rotate=0]
    \draw[sheet] (-1.5,-0.65) rectangle (1.5,0.65);
    \node[mainLabel] at (0,0.95) {$C_p\!\left(K^{(2)}\right)$};

    % miniature simplex icons / coordinates
    \coordinate (s2a) at (-0.95,-0.05);
    \coordinate (s2b) at (-0.35,0.24);
    \coordinate (s2c) at (0.15,-0.08);
    \coordinate (s2d) at (0.72,0.20);

    \draw[simplexEdge] (s2a)--(s2b)--(s2c);
    \draw[simplexEdge] (s2a)--(s2c);
    \draw[simplexEdge] (s2c)--(s2d);

    \node[simplexNode] at (s2a) {};
    \node[simplexNode] at (s2b) {};
    \node[simplexNode] at (s2c) {};
    \node[simplexNode] at (s2d) {};

    \node[coordLabel] at (-0.95,-0.42) {$e_1$};
    \node[coordLabel] at (-0.12,-0.42) {$e_3$};
    \node[coordLabel] at (0.72,-0.42) {$e_4$};
\end{scope}

% ----------------------------
% Registration gate
% ----------------------------
\draw[gate] (5.45,0.75) rectangle (7.10,4.15);

% crosshair corners
\draw[framegray] (5.62,3.98) -- (5.78,3.98);
\draw[framegray] (5.62,3.98) -- (5.62,3.82);

\draw[framegray] (6.93,3.98) -- (6.77,3.98);
\draw[framegray] (6.93,3.98) -- (6.93,3.82);

\draw[framegray] (5.62,0.92) -- (5.78,0.92);
\draw[framegray] (5.62,0.92) -- (5.62,1.08);

\draw[framegray] (6.93,0.92) -- (6.77,0.92);
\draw[framegray] (6.93,0.92) -- (6.93,1.08);

% faint internal grid
\draw[ghostgray!70, line width=0.4pt] (6.28,0.95) -- (6.28,3.95);
\draw[ghostgray!70, line width=0.4pt] (5.70,2.45) -- (6.86,2.45);

\node[smallLabel, align=center, fill=lightfill, inner sep=1pt] at (6.28,2.72) {aligned\\reference set};
\node[mainLabel, fill=lightfill, inner sep=1pt] at (6.28,1.90) {$\Sigma_p$};

% ----------------------------
% Mapping arrows
% ----------------------------
\draw[mapArrow] (3.90,3.25) to[out=0,in=180] node[smallLabel, above] {$E_p^{(1)}$} (5.55,3.05);
\draw[mapArrow] (3.90,1.65) to[out=0,in=180] node[smallLabel, below] {$E_p^{(2)}$} (5.55,1.85);

% ----------------------------
% Ambient panel
% ----------------------------
\draw[ambient] (8.15,0.60) rectangle (14.25,4.35);
\node[mainLabel] at (11.20,4.65) {$\widetilde{C}_p \cong \mathbb{R}^{N_p}$};

% separators
\draw[ghostgray, line width=0.6pt] (10.00,0.85) -- (10.00,4.10);
\draw[ghostgray, line width=0.6pt] (11.80,0.85) -- (11.80,4.10);

% headers
\node[smallLabel] at (9.18,3.98) {basis};
\node[smallLabel] at (10.82,3.98) {$x^{(1)}$};
\node[smallLabel] at (12.42,3.98) {$x^{(2)}$};

% basis coordinates
\node[coordLabel] at (9.15,3.45) {$e_1$};
\node[coordLabel] at (9.15,2.90) {$e_2$};
\node[coordLabel] at (9.15,2.35) {$e_3$};
\node[coordLabel] at (9.15,1.80) {$e_4$};
\node[coordLabel] at (9.15,1.25) {$\vdots$};

% Slots for x^(1)
\draw[solidSlot] (10.45,3.35) rectangle (11.20,3.55);
\draw[solidSlot] (10.45,2.80) rectangle (11.20,3.00);
\draw[zeroSlot]  (10.45,2.25) rectangle (11.20,2.45);
\draw[solidSlot] (10.45,1.70) rectangle (11.20,1.90);
\draw[zeroSlot]  (10.45,1.15) rectangle (11.20,1.35);

\node[smallLabel] at (10.82,3.45) {$\hat c^{(1)}_{1}$};
\node[smallLabel] at (10.82,2.90) {$\hat c^{(1)}_{2}$};
\node[zeroLabel]  at (10.82,2.35) {$0$};
\node[smallLabel] at (10.82,1.80) {$\hat c^{(1)}_{4}$};
\node[zeroLabel]  at (10.82,1.25) {$\cdots$};

% Slots for x^(2)
\draw[solidSlot] (12.05,3.35) rectangle (12.80,3.55);
\draw[zeroSlot]  (12.05,2.80) rectangle (12.80,3.00);
\draw[solidSlot] (12.05,2.25) rectangle (12.80,2.45);
\draw[solidSlot] (12.05,1.70) rectangle (12.80,1.90);
\draw[zeroSlot]  (12.05,1.15) rectangle (12.80,1.35);

\node[smallLabel] at (12.42,3.45) {$\hat c^{(2)}_{1}$};
\node[zeroLabel]  at (12.42,2.90) {$0$};
\node[smallLabel] at (12.42,2.35) {$\hat c^{(2)}_{3}$};
\node[smallLabel] at (12.42,1.80) {$\hat c^{(2)}_{4}$};
\node[zeroLabel]  at (12.42,1.25) {$\cdots$};

% right-side note
\node[smallLabel, align=left, text=framegray, anchor=west, fill=white, inner sep=1pt] at (13.10,1.95) {absent simplex\\$\rightarrow 0$};

% ----------------------------
% Rightward arrow from gate to ambient space
% ----------------------------
\draw[mapArrow] (7.10,2.45) -- (8.15,2.45);

\end{tikzpicture}
\caption{
Alignment of layerwise chain spaces into a shared ambient basis. Each layer is embedded by a coordinate-alignment map $E_p^{(i)}$ into the common ambient chain space $\widetilde{C}_p$, with missing simplices represented by zero-padding so that cross-layer coupling is defined on matched coordinates.
}
\label{fig:alignment}
\end{figure*}
}

% ============================================================================
% Figure A4
% Name: Block structure of the Supra-Hodge Laplacian
% Purpose: Matrix-level diagram separating diagonal Hodge terms and off-diagonal
%          coupling terms.
% Insert near: Definition of L_p (immediately after block definition)
% ============================================================================
\newcommand{\FigureAFour}{
\begin{figure*}[t]
\centering
\begin{tikzpicture}[x=1cm,y=1cm,>=Latex,scale=1.12,every node/.style={transform shape}]

% ----------------------------
% Color palette
% ----------------------------
\definecolor{topoblue}{RGB}{70,120,200}
\definecolor{couplered}{RGB}{190,95,95}
\definecolor{harmgreen}{RGB}{80,150,110}
\definecolor{ghostgray}{RGB}{180,180,180}
\definecolor{framegray}{RGB}{80,80,80}
\definecolor{lightfill}{RGB}{245,247,250}

% ----------------------------
% Styles
% ----------------------------
\tikzset{
    outerFrame/.style={
        rounded corners=4pt,
        draw=framegray,
        line width=0.85pt
    },
    diagBlock/.style={
        rounded corners=2pt,
        draw=topoblue!85!black,
        fill=topoblue!10,
        line width=0.9pt
    },
    offBlock/.style={
        rounded corners=2pt,
        draw=couplered!85!black,
        fill=couplered!8,
        line width=0.85pt
    },
    insetFrame/.style={
        rounded corners=3pt,
        draw=framegray,
        fill=white,
        line width=0.8pt
    },
    cable/.style={->, draw=framegray, line width=0.95pt},
    outputArrow/.style={->, draw=framegray, line width=1.05pt},
    swirl/.style={->, draw=topoblue!80, line width=0.8pt},
    titleStyle/.style={font=\small},
    miniStyle/.style={font=\scriptsize},
    noteStyle/.style={font=\scriptsize, text=framegray},
    mathStyle/.style={font=\scriptsize},
    dotsStyle/.style={font=\large, text=ghostgray}
}

% ----------------------------
% Title
% ----------------------------
\node[titleStyle] at (6.3,5.40) {Supra-Hodge assembly};

% ----------------------------
% Left input signals
% ----------------------------
\draw[cable] (0.7,4.0) -- (2.3,4.0);
\draw[cable] (0.7,2.7) -- (2.3,2.7);
\draw[cable] (0.7,1.4) -- (2.3,1.4);

\node[miniStyle, anchor=east] at (0.55,4.0) {$v^{(1)}$};
\node[miniStyle, anchor=east] at (0.55,2.7) {$v^{(2)}$};
\node[miniStyle, anchor=east] at (0.55,1.4) {$v^{(3)}$};

% ----------------------------
% Central machine frame
% ----------------------------
\draw[outerFrame, fill=lightfill] (3.0,0.8) rectangle (9.7,4.6);

% faint matrix guides
\draw[ghostgray!70, line width=0.45pt] (5.3,0.95) -- (5.3,4.45);
\draw[ghostgray!70, line width=0.45pt] (7.4,0.95) -- (7.4,4.45);
\draw[ghostgray!70, line width=0.45pt] (3.15,3.35) -- (9.55,3.35);
\draw[ghostgray!70, line width=0.45pt] (3.15,2.05) -- (9.55,2.05);

% ----------------------------
% Diagonal Hodge blocks
% ----------------------------
\draw[diagBlock] (3.45,3.65) rectangle (5.05,4.35);
\draw[diagBlock] (5.55,2.35) rectangle (7.15,3.05);
\draw[diagBlock] (7.65,1.05) rectangle (9.25,1.75);

\node[mathStyle] at (4.25,4.00) {$H_p^{(1)} + d_1 I$};
\node[mathStyle] at (6.35,2.70) {$H_p^{(2)} + d_2 I$};
\node[mathStyle] at (8.45,1.40) {$H_p^{(3)} + d_3 I$};

% tiny internal swirl cues
\draw[swirl] (4.55,4.00) arc[start angle=30,end angle=320,radius=0.11];
\draw[swirl] (6.65,2.70) arc[start angle=30,end angle=320,radius=0.11];
\draw[swirl] (8.75,1.40) arc[start angle=30,end angle=320,radius=0.11];

% ----------------------------
% Off-diagonal coupling blocks
% ----------------------------
\draw[offBlock] (5.72,3.77) rectangle (6.98,4.23);
\draw[offBlock] (3.62,2.47) rectangle (4.88,2.93);
\draw[offBlock] (7.82,2.47) rectangle (9.08,2.93);
\draw[offBlock] (5.72,1.17) rectangle (6.98,1.63);

\node[mathStyle] at (6.35,4.00) {$-\omega_{12} I$};
\node[mathStyle] at (4.25,2.70) {$-\omega_{21} I$};
\node[mathStyle] at (8.45,2.70) {$-\omega_{23} I$};
\node[mathStyle] at (6.35,1.40) {$-\omega_{32} I$};

% faded omitted entries
\node[dotsStyle] at (8.45,4.00) {$\cdots$};
\node[dotsStyle] at (4.25,1.40) {$\cdots$};

% ----------------------------
% Input arrows into machine
% ----------------------------
\draw[cable] (2.3,4.0) -- (3.45,4.0);
\draw[cable] (2.3,2.7) -- (3.62,2.7);
\draw[cable] (2.3,1.4) -- (5.72,1.4);

% ----------------------------
% Output arrow to matrix inset
% ----------------------------
\draw[outputArrow] (9.7,2.7) -- (10.5,2.7);

% ----------------------------
% Matrix inset
% ----------------------------
\draw[insetFrame] (10.60,1.15) rectangle (13.85,4.25);

\node[titleStyle] at (12.22,4.55) {$L_p$};

\node[mathStyle, align=center] at (12.22,2.70) {%
$\displaystyle
L_p=
\begin{bmatrix}
H_p^{(1)}+d_1I & -\omega_{12}I & \cdots \\
-\omega_{21}I & H_p^{(2)}+d_2I & \cdots \\
\vdots & \vdots & \ddots
\end{bmatrix}
$};

% ----------------------------
% Bottom notes
% ----------------------------
\node[noteStyle, align=center, fill=white, inner sep=1pt] at (4.05,0.20) {diagonal: intra-layer topology};
\node[noteStyle, align=center, fill=white, inner sep=1pt] at (8.95,0.20) {off-diagonal: cross-layer coupling};

% ----------------------------
% Optional labels inside main frame
% ----------------------------
\node[miniStyle, text=framegray, fill=lightfill, inner sep=1pt] at (4.55,4.42) {block operator view};
\node[miniStyle, text=framegray, fill=lightfill, inner sep=1pt] at (8.10,4.42) {coupled assembly};

\end{tikzpicture}
\caption{
Block-operator view of the Supra-Hodge assembly. Diagonal blocks encode layerwise Hodge structure augmented by coupling-degree terms, while off-diagonal blocks encode inter-layer attraction. The machine-like layout preserves the matrix semantics of $L_p$ while making the separation between intra-layer topology and cross-layer coupling visually explicit.
}
\label{fig:block}
\end{figure*}
}

% ============================================================================
% Figure A5
% Name: Energy decomposition
% Purpose: Geometry-meets-algebra picture for v^T L_p v decomposition.
% Insert near: Fundamental properties section, immediately after Prop. sh_energy
% ============================================================================
\newcommand{\FigureAFive}{
\begin{figure*}[t]
\centering
\begin{tikzpicture}[x=1cm,y=1cm,>=Latex,scale=1.12,every node/.style={transform shape}]

% ----------------------------
% Color palette
% ----------------------------
\definecolor{topoblue}{RGB}{70,120,200}
\definecolor{couplered}{RGB}{190,95,95}
\definecolor{harmgreen}{RGB}{80,150,110}
\definecolor{softgold}{RGB}{210,170,80}
\definecolor{ghostgray}{RGB}{180,180,180}
\definecolor{framegray}{RGB}{80,80,80}
\definecolor{lightfill}{RGB}{245,247,250}

% ----------------------------
% Styles
% ----------------------------
\tikzset{
    sourcePanel/.style={
        rounded corners=4pt,
        draw=framegray,
        fill=lightfill,
        line width=0.8pt
    },
    topoEdge/.style={draw=topoblue, line width=0.95pt},
    faintEdge/.style={draw=ghostgray, line width=0.7pt},
    coupling/.style={draw=couplered!80, line width=0.85pt},
    softCoupling/.style={draw=couplered!60, dashed, line width=0.75pt},
    basinLine/.style={draw=framegray, line width=1.0pt},
    flowArrow/.style={->, draw=framegray, line width=0.95pt},
    smallArrow/.style={->, draw=topoblue, line width=0.85pt},
    titleStyle/.style={font=\small},
    miniStyle/.style={font=\scriptsize},
    noteStyle/.style={font=\scriptsize, text=framegray},
    mathStyle/.style={font=\small}
}

% ----------------------------
% Central equation
% ----------------------------
\draw[sourcePanel, fill=white] (5.82,4.70) rectangle (8.18,5.18);
\node[mathStyle] at (7.0,4.94) {$\displaystyle v^\top L_p v$};

% ----------------------------
% Left source panel: within-layer roughness
% ----------------------------
\draw[sourcePanel] (1.0,2.5) rectangle (5.0,4.5);
\node[titleStyle] at (3.0,4.78) {within-layer roughness};

% simplicial layer
\coordinate (La) at (1.8,3.0);
\coordinate (Lb) at (3.0,2.8);
\coordinate (Lc) at (4.1,3.3);
\coordinate (Ld) at (3.3,4.0);
\coordinate (Le) at (2.1,4.1);

% faces
\fill[topoblue!16] (La)--(Lb)--(Ld)--cycle;
\fill[topoblue!10] (La)--(Ld)--(Le)--cycle;

% edges
\draw[topoEdge] (La)--(Lb)--(Lc);
\draw[topoEdge] (La)--(Le)--(Ld)--(Lc);
\draw[topoEdge] (Lb)--(Ld);
\draw[topoEdge] (La)--(Ld);

% nodes
\fill[white] (La) circle (0.06);
\fill[white] (Lb) circle (0.06);
\fill[white] (Lc) circle (0.06);
\fill[white] (Ld) circle (0.06);
\fill[white] (Le) circle (0.06);

\draw[framegray, line width=0.7pt] (La) circle (0.06);
\draw[framegray, line width=0.7pt] (Lb) circle (0.06);
\draw[framegray, line width=0.7pt] (Lc) circle (0.06);
\draw[framegray, line width=0.7pt] (Ld) circle (0.06);
\draw[framegray, line width=0.7pt] (Le) circle (0.06);

% roughness cues
\draw[smallArrow] (La) -- ($(Ld)!0.82!(La)$);
\fill[couplered!35] (Lb) circle (0.10);

% faint contour / roughness field
\draw[ghostgray!80, line width=0.5pt]
    (1.55,3.55) .. controls (2.15,3.85) and (3.05,3.15) .. (4.30,3.65);
\draw[ghostgray!60, line width=0.45pt]
    (1.55,3.30) .. controls (2.25,3.55) and (3.15,2.95) .. (4.20,3.40);

\node[miniStyle] at (3.0,2.18)
{$\displaystyle \sum_i \langle v^{(i)}, H_p^{(i)} v^{(i)} \rangle$};

% ----------------------------
% Right source panel: cross-layer disagreement
% ----------------------------
\draw[sourcePanel] (9.0,2.5) rectangle (13.0,4.5);
\node[titleStyle] at (11.0,4.78) {cross-layer disagreement};

% upper layer
\coordinate (Ua) at (10.2,3.65);
\coordinate (Ub) at (11.8,3.65);
\coordinate (Uc) at (10.2,4.25);
\coordinate (Ud) at (11.8,4.25);

% lower layer
\coordinate (Va) at (10.2,2.85);
\coordinate (Vb) at (11.8,2.85);
\coordinate (Vc) at (10.2,3.45);
\coordinate (Vd) at (11.8,3.45);

% upper layer face and edges
\fill[topoblue!16] (Ua)--(Ub)--(Ud)--cycle;
\draw[topoEdge] (Ua)--(Ub)--(Ud)--(Uc)--cycle;
\draw[topoEdge] (Ua)--(Ud);

% lower layer face and edges
\fill[couplered!14] (Va)--(Vc)--(Vd)--cycle;
\draw[couplered!80, line width=0.95pt] (Va)--(Vb)--(Vd)--(Vc)--cycle;
\draw[couplered!80, line width=0.95pt] (Va)--(Vd);

% nodes upper
\foreach \pt in {Ua,Ub,Uc,Ud} {
    \fill[white] (\pt) circle (0.055);
    \draw[framegray, line width=0.7pt] (\pt) circle (0.055);
}
% nodes lower
\foreach \pt in {Va,Vb,Vc,Vd} {
    \fill[white] (\pt) circle (0.055);
    \draw[framegray, line width=0.7pt] (\pt) circle (0.055);
}

% vertical couplings
\draw[softCoupling] (Ua)--(Va);
\draw[softCoupling] (Ub)--(Vb);
\draw[softCoupling] (Uc)--(Vc);
\draw[softCoupling] (Ud)--(Vd);

\node[miniStyle] at (11.0,2.18)
{$\displaystyle \sum_{i<j}\omega_{ij}\|v^{(i)}-v^{(j)}\|^2$};

% ----------------------------
% Tributary arrows into basin
% ----------------------------
\draw[flowArrow] (5.0,3.20) .. controls (5.85,2.85) and (6.20,2.25) .. (6.92,1.80);
\draw[flowArrow] (9.0,3.20) .. controls (8.15,2.85) and (7.80,2.25) .. (7.08,1.80);

% ----------------------------
% Merge basin
% ----------------------------
% basin rim
\draw[basinLine]
    (5.70,1.55)
    .. controls (6.10,1.25) and (6.55,1.05) .. (7.00,1.05)
    .. controls (7.45,1.05) and (7.90,1.25) .. (8.30,1.55);

% basin shoulders
\draw[basinLine] (5.35,1.82) .. controls (5.50,1.68) and (5.60,1.60) .. (5.70,1.55);
\draw[basinLine] (8.65,1.82) .. controls (8.50,1.68) and (8.40,1.60) .. (8.30,1.55);

% light fill in basin
\fill[lightfill]
    (5.70,1.55)
    .. controls (6.10,1.25) and (6.55,1.05) .. (7.00,1.05)
    .. controls (7.45,1.05) and (7.90,1.25) .. (8.30,1.55)
    -- (8.30,1.62) -- (5.70,1.62) -- cycle;

\node[miniStyle, fill=lightfill, inner sep=1pt] at (7.0,1.35) {total energy};

% ----------------------------
% Tiny visual hints beneath equation
% ----------------------------
\node[noteStyle, fill=white, inner sep=1pt] at (3.0,0.82) {local topological incompatibility};
\node[noteStyle, fill=white, inner sep=1pt] at (11.0,0.82) {cross-layer mismatch};

\end{tikzpicture}
\caption{
Energy decomposition for the Supra-Hodge operator. The total quadratic form $v^\top L_p v$ splits into a within-layer contribution measuring topological roughness inside each layer and an inter-layer contribution measuring disagreement across coupled layers. The tributary layout emphasizes that both sources feed the same total energy landscape.
}
\label{fig:energy}
\end{figure*}
}

% ============================================================================
% Figure A6
% Name: Consensus-harmonic nullspace
% Purpose: Visual intuition for theorem: harmonicity in each layer + agreement
%          across coupled layers.
% Insert near: Nullspace intuition / just before main theorem
% ============================================================================
\newcommand{\FigureASix}{
\begin{figure*}[t]
\centering
\begin{tikzpicture}[x=1cm,y=1cm,>=Latex,scale=1.10,every node/.style={transform shape}]

% ----------------------------
% Color palette
% ----------------------------
\definecolor{topoblue}{RGB}{70,120,200}
\definecolor{couplered}{RGB}{190,95,95}
\definecolor{harmgreen}{RGB}{80,150,110}
\definecolor{ghostgray}{RGB}{180,180,180}
\definecolor{framegray}{RGB}{80,80,80}
\definecolor{lightfill}{RGB}{245,247,250}

% ----------------------------
% Styles
% ----------------------------
\tikzset{
    panelFrame/.style={
        rounded corners=4pt,
        draw=framegray,
        fill=lightfill,
        line width=0.8pt
    },
    layerEdge/.style={draw=framegray, line width=0.9pt},
    goodEdge/.style={draw=harmgreen!85!black, line width=0.95pt},
    badEdge/.style={draw=couplered!85!black, line width=0.95pt},
    alignGood/.style={draw=harmgreen!75!black, line width=0.8pt},
    alignBad/.style={draw=couplered!80, line width=0.8pt},
    nodeStyle/.style={
        circle, draw=framegray, fill=white,
        inner sep=0pt, minimum size=4.6pt, line width=0.75pt
    },
    titleStyle/.style={font=\small},
    miniStyle/.style={font=\scriptsize},
    noteStyle/.style={font=\scriptsize, text=framegray},
    mathStyle/.style={font=\scriptsize}
}

% ----------------------------
% Panel frames
% ----------------------------
\draw[panelFrame] (0.5,0.75) rectangle (6.2,5.00);
\draw[panelFrame] (7.0,0.75) rectangle (12.7,5.00);

\node[titleStyle] at (3.35,5.32) {consensus-harmonic mode};
\node[titleStyle] at (9.85,5.32) {non-consensus contrast};

% divider note
\node[miniStyle, text=framegray, fill=white, inner sep=1pt] at (6.60,5.12) {$\ker(L_p)$};

% ----------------------------
% LEFT PANEL: consensus-harmonic
% ----------------------------

% layer centers
\coordinate (L1) at (3.35,4.10);
\coordinate (L2) at (3.35,2.90);
\coordinate (L3) at (3.35,1.70);

% --- Layer 1
\coordinate (l1a) at ($(L1)+(-0.75,-0.25)$);
\coordinate (l1b) at ($(L1)+( 0.75,-0.25)$);
\coordinate (l1c) at ($(L1)+(-0.75, 0.45)$);
\coordinate (l1d) at ($(L1)+( 0.75, 0.45)$);

\fill[harmgreen!16] (l1a)--(l1b)--(l1d)--cycle;
\draw[goodEdge] (l1a)--(l1b)--(l1d)--(l1c)--cycle;
\draw[goodEdge] (l1a)--(l1d);
\draw[->, draw=harmgreen!85!black, line width=0.85pt, shorten <=2pt, shorten >=2pt] (l1a)--(l1d);

\node[nodeStyle] at (l1a) {};
\node[nodeStyle] at (l1b) {};
\node[nodeStyle] at (l1c) {};
\node[nodeStyle] at (l1d) {};

% --- Layer 2
\coordinate (l2a) at ($(L2)+(-0.75,-0.25)$);
\coordinate (l2b) at ($(L2)+( 0.75,-0.25)$);
\coordinate (l2c) at ($(L2)+(-0.75, 0.45)$);
\coordinate (l2d) at ($(L2)+( 0.75, 0.45)$);

\fill[harmgreen!16] (l2a)--(l2b)--(l2d)--cycle;
\draw[goodEdge] (l2a)--(l2b)--(l2d)--(l2c)--cycle;
\draw[goodEdge] (l2a)--(l2d);
\draw[->, draw=harmgreen!85!black, line width=0.85pt, shorten <=2pt, shorten >=2pt] (l2a)--(l2d);

\node[nodeStyle] at (l2a) {};
\node[nodeStyle] at (l2b) {};
\node[nodeStyle] at (l2c) {};
\node[nodeStyle] at (l2d) {};

% --- Layer 3
\coordinate (l3a) at ($(L3)+(-0.75,-0.25)$);
\coordinate (l3b) at ($(L3)+( 0.75,-0.25)$);
\coordinate (l3c) at ($(L3)+(-0.75, 0.45)$);
\coordinate (l3d) at ($(L3)+( 0.75, 0.45)$);

\fill[harmgreen!16] (l3a)--(l3b)--(l3d)--cycle;
\draw[goodEdge] (l3a)--(l3b)--(l3d)--(l3c)--cycle;
\draw[goodEdge] (l3a)--(l3d);
\draw[->, draw=harmgreen!85!black, line width=0.85pt, shorten <=2pt, shorten >=2pt] (l3a)--(l3d);

\node[nodeStyle] at (l3a) {};
\node[nodeStyle] at (l3b) {};
\node[nodeStyle] at (l3c) {};
\node[nodeStyle] at (l3d) {};

% vertical alignment links
\draw[alignGood] (l1a)--(l2a)--(l3a);
\draw[alignGood] (l1b)--(l2b)--(l3b);
\draw[alignGood] (l1c)--(l2c)--(l3c);
\draw[alignGood] (l1d)--(l2d)--(l3d);

% smooth shared wave trace
\draw[harmgreen!70, line width=0.8pt]
    (1.25,4.28)
    .. controls (2.05,4.42) and (2.65,3.95) .. (3.35,4.12)
    .. controls (4.05,4.28) and (4.55,3.82) .. (5.15,3.98)
    .. controls (5.00,3.55) and (4.60,3.35) .. (4.10,3.18)
    .. controls (3.55,3.00) and (3.00,2.72) .. (2.55,2.55)
    .. controls (2.10,2.38) and (1.75,2.10) .. (1.45,1.96)
    .. controls (2.05,2.12) and (2.60,1.72) .. (3.30,1.88)
    .. controls (4.00,2.05) and (4.55,1.55) .. (5.10,1.72);

% equations
\node[mathStyle, fill=lightfill, inner sep=1pt] at (5.00,4.10) {$H_p^{(1)}v^{(1)}=0$};
\node[mathStyle, fill=lightfill, inner sep=1pt] at (5.00,2.90) {$H_p^{(2)}v^{(2)}=0$};
\node[mathStyle, fill=lightfill, inner sep=1pt] at (5.00,1.70) {$H_p^{(3)}v^{(3)}=0$};

\node[miniStyle] at (3.35,0.42) {$v^{(1)} = v^{(2)} = v^{(3)}$};

% ----------------------------
% RIGHT PANEL: non-consensus contrast
% ----------------------------

% layer centers
\coordinate (R1) at (9.65,4.10);
\coordinate (R2) at (9.65,2.90);
\coordinate (R3) at (9.65,1.70);

% --- Top layer
\coordinate (r1a) at ($(R1)+(-0.75,-0.25)$);
\coordinate (r1b) at ($(R1)+( 0.75,-0.25)$);
\coordinate (r1c) at ($(R1)+(-0.75, 0.45)$);
\coordinate (r1d) at ($(R1)+( 0.75, 0.45)$);

\fill[topoblue!16] (r1a)--(r1b)--(r1d)--cycle;
\draw[layerEdge] (r1a)--(r1b)--(r1d)--(r1c)--cycle;
\draw[layerEdge] (r1a)--(r1d);
\draw[->, draw=topoblue!90!black, line width=0.85pt, shorten <=2pt, shorten >=2pt] (r1a)--(r1d);

\node[nodeStyle] at (r1a) {};
\node[nodeStyle] at (r1b) {};
\node[nodeStyle] at (r1c) {};
\node[nodeStyle] at (r1d) {};

% --- Middle layer (flipped / discordant)
\coordinate (r2a) at ($(R2)+(-0.75,-0.25)$);
\coordinate (r2b) at ($(R2)+( 0.75,-0.25)$);
\coordinate (r2c) at ($(R2)+(-0.75, 0.45)$);
\coordinate (r2d) at ($(R2)+( 0.75, 0.45)$);

\fill[couplered!16] (r2a)--(r2b)--(r2d)--cycle;
\draw[badEdge] (r2a)--(r2b)--(r2d)--(r2c)--cycle;
\draw[badEdge] (r2a)--(r2d);
\draw[->, draw=couplered!90!black, line width=0.85pt, shorten <=2pt, shorten >=2pt] (r2d)--(r2a);

\node[nodeStyle] at (r2a) {};
\node[nodeStyle] at (r2b) {};
\node[nodeStyle] at (r2c) {};
\node[nodeStyle] at (r2d) {};

% --- Bottom layer
\coordinate (r3a) at ($(R3)+(-0.75,-0.25)$);
\coordinate (r3b) at ($(R3)+( 0.75,-0.25)$);
\coordinate (r3c) at ($(R3)+(-0.75, 0.45)$);
\coordinate (r3d) at ($(R3)+( 0.75, 0.45)$);

\fill[topoblue!16] (r3a)--(r3b)--(r3d)--cycle;
\draw[layerEdge] (r3a)--(r3b)--(r3d)--(r3c)--cycle;
\draw[layerEdge] (r3a)--(r3d);
\draw[->, draw=topoblue!90!black, line width=0.85pt, shorten <=2pt, shorten >=2pt] (r3a)--(r3d);

\node[nodeStyle] at (r3a) {};
\node[nodeStyle] at (r3b) {};
\node[nodeStyle] at (r3c) {};
\node[nodeStyle] at (r3d) {};

% broken / zig-zag couplings
\draw[alignBad] (r1a) -- ($(r1a)!0.45!(r2a)+(0.08,-0.02)$) -- (r2a);
\draw[alignBad] (r1b) -- ($(r1b)!0.45!(r2b)+(-0.08,-0.02)$) -- (r2b);
\draw[alignBad] (r1c) -- ($(r1c)!0.45!(r2c)+(0.08,0.00)$) -- (r2c);
\draw[alignBad] (r1d) -- ($(r1d)!0.45!(r2d)+(-0.08,0.00)$) -- (r2d);

\draw[alignBad] (r2a) -- ($(r2a)!0.45!(r3a)+(-0.08,0.02)$) -- (r3a);
\draw[alignBad] (r2b) -- ($(r2b)!0.45!(r3b)+(0.08,0.02)$) -- (r3b);
\draw[alignBad] (r2c) -- ($(r2c)!0.45!(r3c)+(-0.08,0.00)$) -- (r3c);
\draw[alignBad] (r2d) -- ($(r2d)!0.45!(r3d)+(0.08,0.00)$) -- (r3d);

\node[miniStyle] at (9.85,0.42) {agreement broken across layers};

\end{tikzpicture}
\caption{
Nullspace intuition for the Supra-Hodge operator. On the left, a stacked signal lies in $\ker(L_p)$ because it is harmonic in every layer and aligned across coupled layers. On the right, a cross-layer contrast breaks consensus, producing disagreement even when individual layer structures remain visually similar.
}
\label{fig:nullspace}
\end{figure*}
}

% ============================================================================
% Figure A7
% Name: Spectral modes
% Purpose: Show qualitative eigenmode regimes (near-zero, Fiedler-like, high).
% Insert near: Spectral interpretation section introduction
% ============================================================================
\newcommand{\FigureASeven}{
\begin{figure*}[t]
\centering
\begin{tikzpicture}[x=1cm,y=1cm,>=Latex,scale=1.10,every node/.style={transform shape}]

% ----------------------------
% Color palette
% ----------------------------
\definecolor{topoblue}{RGB}{70,120,200}
\definecolor{couplered}{RGB}{190,95,95}
\definecolor{harmgreen}{RGB}{80,150,110}
\definecolor{ghostgray}{RGB}{180,180,180}
\definecolor{framegray}{RGB}{80,80,80}
\definecolor{lightfill}{RGB}{245,247,250}

% ----------------------------
% Styles
% ----------------------------
\tikzset{
    panelFrame/.style={
        rounded corners=4pt,
        draw=framegray,
        fill=lightfill,
        line width=0.8pt
    },
    baseEdge/.style={draw=framegray, line width=0.9pt},
    nodeStyle/.style={
        circle, draw=framegray, fill=white,
        inner sep=0pt, minimum size=4.8pt, line width=0.75pt
    },
    titleStyle/.style={font=\small},
    miniStyle/.style={font=\scriptsize},
    noteStyle/.style={font=\scriptsize, text=framegray},
    freqArrow/.style={->, draw=framegray, line width=0.95pt}
}

% ----------------------------
% Panel frames
% ----------------------------
\draw[panelFrame] (0.5,0.8) rectangle (4.2,4.0);
\draw[panelFrame] (5.0,0.8) rectangle (8.7,4.0);
\draw[panelFrame] (9.5,0.8) rectangle (13.2,4.0);

% ----------------------------
% Panel 1: near-zero coherent mode
% ----------------------------
\node[titleStyle] at (2.35,4.25) {$\lambda \approx 0$};

% calm background field
\fill[harmgreen!10] (2.35,2.40) ellipse (1.45 and 1.05);
\fill[topoblue!8]  (2.20,2.35) ellipse (1.10 and 0.82);

% graph coordinates
\coordinate (p1a) at (1.40,2.00);
\coordinate (p1b) at (2.35,1.65);
\coordinate (p1c) at (3.25,2.20);
\coordinate (p1d) at (2.75,3.15);
\coordinate (p1e) at (1.55,3.05);

\draw[baseEdge] (p1a)--(p1b)--(p1c)--(p1d)--(p1e)--cycle;
\draw[baseEdge] (p1a)--(p1d);
\draw[baseEdge] (p1b)--(p1d);

% nearly uniform node tones
\fill[harmgreen!20] (p1a) circle (0.11);
\fill[harmgreen!18] (p1b) circle (0.11);
\fill[harmgreen!22] (p1c) circle (0.11);
\fill[harmgreen!17] (p1d) circle (0.11);
\fill[harmgreen!19] (p1e) circle (0.11);

\node[nodeStyle] at (p1a) {};
\node[nodeStyle] at (p1b) {};
\node[nodeStyle] at (p1c) {};
\node[nodeStyle] at (p1d) {};
\node[nodeStyle] at (p1e) {};

\node[miniStyle] at (2.35,0.68) {coherent plateau};

% ----------------------------
% Panel 2: Fiedler-type split
% ----------------------------
\node[titleStyle] at (6.85,4.25) {$\lambda_2$};

% split background
\fill[topoblue!12] (5.15,0.95) rectangle (6.95,3.85);
\fill[couplered!12] (6.95,0.95) rectangle (8.55,3.85);

% interface / fault line
\draw[ghostgray!90, dashed, line width=0.8pt]
    (6.95,1.10) .. controls (7.10,1.95) and (6.85,2.55) .. (7.00,3.55);

% graph coordinates
\coordinate (p2a) at (5.90,2.00);
\coordinate (p2b) at (6.85,1.65);
\coordinate (p2c) at (7.75,2.20);
\coordinate (p2d) at (7.25,3.15);
\coordinate (p2e) at (6.05,3.05);

\draw[baseEdge] (p2a)--(p2b)--(p2c)--(p2d)--(p2e)--cycle;
\draw[baseEdge] (p2a)--(p2d);
\draw[baseEdge] (p2b)--(p2d);

% signed partition node tones
\fill[topoblue!28]   (p2a) circle (0.11);
\fill[ghostgray!35]  (p2b) circle (0.11);
\fill[couplered!28]  (p2c) circle (0.11);
\fill[couplered!22]  (p2d) circle (0.11);
\fill[topoblue!22]   (p2e) circle (0.11);

\node[nodeStyle] at (p2a) {};
\node[nodeStyle] at (p2b) {};
\node[nodeStyle] at (p2c) {};
\node[nodeStyle] at (p2d) {};
\node[nodeStyle] at (p2e) {};

\node[miniStyle] at (6.85,0.68) {signed partition};

% ----------------------------
% Panel 3: high-frequency oscillatory mode
% ----------------------------
\node[titleStyle] at (11.35,4.25) {$\lambda$ large};

% alternating diagonal bands
\begin{scope}
\clip (9.6,0.9) rectangle (13.1,3.9);
\fill[topoblue!10]  (8.9,1.0) -- (9.8,0.9) -- (12.7,3.9) -- (11.8,4.0) -- cycle;
\fill[couplered!10] (9.4,0.9) -- (10.3,0.9) -- (13.2,3.9) -- (12.3,3.9) -- cycle;
\fill[topoblue!10]  (9.9,0.9) -- (10.8,0.9) -- (13.7,3.9) -- (12.8,3.9) -- cycle;
\fill[couplered!10] (10.4,0.9) -- (11.3,0.9) -- (14.2,3.9) -- (13.3,3.9) -- cycle;
\end{scope}

% graph coordinates
\coordinate (p3a) at (10.40,2.00);
\coordinate (p3b) at (11.35,1.65);
\coordinate (p3c) at (12.25,2.20);
\coordinate (p3d) at (11.75,3.15);
\coordinate (p3e) at (10.55,3.05);

\draw[baseEdge] (p3a)--(p3b)--(p3c)--(p3d)--(p3e)--cycle;
\draw[baseEdge] (p3a)--(p3d);
\draw[baseEdge] (p3b)--(p3d);

% alternating node tones
\fill[topoblue!30]  (p3a) circle (0.11);
\fill[couplered!30] (p3b) circle (0.11);
\fill[topoblue!30]  (p3c) circle (0.11);
\fill[couplered!30] (p3d) circle (0.11);
\fill[topoblue!26]  (p3e) circle (0.11);

\node[nodeStyle] at (p3a) {};
\node[nodeStyle] at (p3b) {};
\node[nodeStyle] at (p3c) {};
\node[nodeStyle] at (p3d) {};
\node[nodeStyle] at (p3e) {};

% optional short opposing arrows
\draw[->, draw=topoblue!85!black, line width=0.75pt, shorten <=2pt, shorten >=2pt] (p3a)--(p3b);
\draw[->, draw=couplered!85!black, line width=0.75pt, shorten <=2pt, shorten >=2pt] (p3d)--(p3b);

\node[miniStyle] at (11.35,0.68) {rapid variation};

% ----------------------------
% Bottom progression arrow
% ----------------------------
\draw[freqArrow] (3.15,-0.02) -- (10.55,-0.02);
\node[noteStyle, fill=white, inner sep=1pt] at (6.85,0.20) {smooth $\rightarrow$ partitioning $\rightarrow$ oscillatory};

\end{tikzpicture}
\caption{
Spectral regimes of the Supra-Hodge operator, shown as progressively more oscillatory mode patterns on the same underlying relational structure. Near-zero modes capture coherent low-energy agreement, intermediate Fiedler-type modes expose structural fault lines or partitions, and high-frequency modes highlight rapidly varying inconsistency patterns.
}
\label{fig:spectrum}
\end{figure*}
}

% ============================================================================
% Figure A8
% Name: End-to-end reasoning pipeline
% Purpose: Three-layer architecture from data to operator construction to
%          spectral diagnostics/outputs.
% Insert near: Application/algorithm pipeline discussion (replaces old pipeline)
% ============================================================================
\newcommand{\FigureAEight}{
\begin{figure*}[t]
\centering
\begin{tikzpicture}[x=1cm,y=1cm,>=Latex,scale=1.12,every node/.style={transform shape}]

% ----------------------------
% Color palette
% ----------------------------
\definecolor{topoblue}{RGB}{70,120,200}
\definecolor{couplered}{RGB}{190,95,95}
\definecolor{harmgreen}{RGB}{80,150,110}
\definecolor{ghostgray}{RGB}{180,180,180}
\definecolor{framegray}{RGB}{80,80,80}
\definecolor{lightfill}{RGB}{245,247,250}

% ----------------------------
% Styles
% ----------------------------
\tikzset{
    room/.style={
        rounded corners=5pt,
        draw=framegray,
        fill=lightfill,
        line width=0.9pt
    },
    panel/.style={
        rounded corners=3pt,
        draw=framegray,
        fill=white,
        line width=0.8pt
    },
    screen/.style={
        rounded corners=2pt,
        draw=framegray,
        fill=topoblue!6,
        line width=0.75pt
    },
    nodeStyle/.style={
        circle, draw=framegray, fill=white,
        inner sep=0pt, minimum size=4.5pt, line width=0.7pt
    },
    edgeStyle/.style={draw=framegray, line width=0.8pt},
    activeEdge/.style={draw=topoblue, line width=0.9pt},
    coupling/.style={draw=couplered!80, dashed, line width=0.8pt},
    arrowMain/.style={->, draw=framegray, line width=1.0pt},
    arrowFlow/.style={->, draw=topoblue, line width=0.9pt},
    titleStyle/.style={font=\small},
    miniStyle/.style={font=\scriptsize},
    noteStyle/.style={font=\scriptsize, text=framegray},
    mathStyle/.style={font=\scriptsize}
}

% ----------------------------
% Outer control room
% ----------------------------
\draw[room] (0.6,0.7) rectangle (13.4,5.2);
\node[titleStyle] at (7.0,5.5) {reasoning control room};

% ----------------------------
% LEFT: Input stack (multi-layer signals)
% ----------------------------
\node[miniStyle] at (2.2,4.9) {layer stack};

\coordinate (A1) at (2.2,3.9);
\coordinate (A2) at (2.2,2.9);
\coordinate (A3) at (2.2,1.9);

% simple layer glyphs
\foreach \c in {A1,A2,A3}{
    \coordinate (n1) at ($(\c)+(-0.6,-0.2)$);
    \coordinate (n2) at ($(\c)+(0.6,-0.2)$);
    \coordinate (n3) at ($(\c)+(0.0,0.5)$);

    \draw[edgeStyle] (n1)--(n2)--(n3)--cycle;
    \node[nodeStyle] at (n1) {};
    \node[nodeStyle] at (n2) {};
    \node[nodeStyle] at (n3) {};
}

% vertical coupling hints
\draw[coupling] (A1) -- (A2) -- (A3);

% ----------------------------
% CENTER: Operator core (L_p)
% ----------------------------
\draw[panel] (4.6,1.4) rectangle (8.4,4.6);
\node[titleStyle] at (6.5,4.9) {$L_p$ operator};

% internal block structure
\draw[screen] (5.0,3.6) rectangle (6.4,4.3);
\draw[screen] (6.6,2.7) rectangle (8.0,3.4);
\draw[screen] (5.0,1.8) rectangle (6.4,2.5);

\node[mathStyle] at (5.7,3.95) {$H^{(1)}$};
\node[mathStyle] at (7.3,3.05) {$H^{(2)}$};
\node[mathStyle] at (5.7,2.15) {$H^{(3)}$};

% coupling blocks
\draw[screen, fill=couplered!8] (6.6,3.6) rectangle (8.0,4.3);
\draw[screen, fill=couplered!8] (5.0,2.7) rectangle (6.4,3.4);

\node[mathStyle] at (7.3,3.95) {$-\omega I$};
\node[mathStyle] at (5.7,3.05) {$-\omega I$};

% arrows from layers into operator
\draw[arrowMain] (2.9,3.9) -- (4.6,3.9);
\draw[arrowMain] (2.9,2.9) -- (4.6,3.0);
\draw[arrowMain] (2.9,1.9) -- (4.6,2.1);

% ----------------------------
% RIGHT: Monitoring screens
% ----------------------------
\node[miniStyle] at (11.0,4.9) {diagnostics};

% screen 1: energy
\draw[panel] (9.4,3.5) rectangle (12.8,4.6);
\node[miniStyle] at (11.1,4.45) {energy};

\draw[topoblue, line width=0.9pt]
    (9.6,3.7) .. controls (10.2,4.2) and (11.2,3.2) .. (12.6,4.2);

% screen 2: spectrum
\draw[panel] (9.4,2.1) rectangle (12.8,3.2);
\node[miniStyle] at (11.1,3.05) {spectrum};

\foreach \x in {9.8,10.4,11.0,11.6,12.2}{
    \draw[framegray, line width=0.7pt] (\x,2.2) -- (\x,2.6);
}

% screen 3: alignment
\draw[panel] (9.4,0.9) rectangle (12.8,2.0);
\node[miniStyle] at (11.1,1.85) {alignment};

\draw[topoblue, line width=0.8pt] (9.7,1.2) -- (12.5,1.7);
\draw[couplered, line width=0.8pt] (9.7,1.7) -- (12.5,1.2);

% ----------------------------
% Output arrow
% ----------------------------
\draw[arrowFlow] (8.4,3.0) -- (9.4,3.0);

\node[noteStyle, fill=white, inner sep=1pt] at (7.0,0.30)
{structured reasoning emerges from coupled higher-order signals};

\end{tikzpicture}
\caption{
A systems-level view of Supra-Hodge reasoning. Multi-layer signals enter a coupled operator $L_p$, which encodes both intra-layer topology and inter-layer alignment. Downstream diagnostics expose energy, spectral structure, and alignment quality, providing interpretable signals for controlling reasoning behavior.
}
\label{fig:pipeline}
\end{figure*}
}
